\section{Open Questions from Paper Re-read}

The following five questions surfaced from a careful re-read of the paper
against the on-disk audit evidence. Each is grounded in a specific
observation from the replication (\S 4 / \S \ref{sec:critique}) and comes
with a concrete, actionable next-step probe.

\begin{enumerate}
  \item \textbf{Does the TOPAS-MGM extension, once released, reproduce the
    paper's He > p MDS-per-dose ratio at matched LET, and is the
    additional ``track-structure'' information relative to the analytical
    MGM core quantifiable as a scalar correction factor per (particle, LET)
    pair?}

    \emph{Basis.} P3 in this audit found MGM He/p ratio $= 0.97$--$1.07$
    at matched LET $\le 35$ keV/$\mu$m, whereas the paper's TOPAS-MGM
    shows He > p. The paper attributes this to denser helium track
    structure, but the analytical MGM has no track-structure knob. Whether
    TOPAS-MGM's excess-over-MGM is a smooth function of LET (correctable
    analytically) or a full-distribution effect (only computable
    per-track) is unresolved and would decide whether analytical MGM can
    be patched or whether the full extension is required.

    \emph{Next steps.} When TOPAS-MGM source is released, run matched-LET
    p vs He at 10 LET points spanning 6--35 keV/$\mu$m; extract
    per-crossing MDS distributions; compute the
    (TOPAS-MGM / analytical-MGM) ratio and check whether it is a smooth
    function of LET. If yes, a simple $1 + \alpha(\text{LET})$ correction
    closes the gap. If no, publish the residual distributions as the true
    ``track-structure signal.''

  \item \textbf{Does the MGM framework, either analytically or via
    TOPAS-MGM, extend cleanly to heavier ions (carbon, oxygen, neon) at
    hadron-therapy LETs (30--200 keV/$\mu$m), or does it break at the
    $b(y_F)$ zero-crossing near $y_F \sim 164$ keV/$\mu$m?}

    \emph{Basis.} The paper explicitly restricts MGM validity to $y_F <
    200$ keV/$\mu$m; our sweep in P2 sees $b(y_F)$ fall to zero near
    164 keV/$\mu$m in the current fits. Carbon at hadron-therapy energies
    routinely exceeds this LET, and mixed-field / heavier-ion FLASH
    regimes are a growing clinical target. Extension to $Z \ge 6$ is not
    tested by the paper and is a natural next-order question.

    \emph{Next steps.} (1) Extract or re-fit $a(y_F), b(y_F)$ at high
    $y_F$ using Geant4-DNA carbon/oxygen radial-energy tables; check
    whether a higher-order polynomial or a different functional form
    removes the $b \to 0$ crossing. (2) Once TOPAS-MGM is available, add
    carbon/oxygen ion source and run against TOPAS-nBio carbon reference
    at 3 LET points. (3) Publish an updated MGM validity range.

  \item \textbf{How does the paper's MDS/complexity output map to
    clonogenic survival, given that MDS is a physics observable and
    survival is the biologically-actionable endpoint?}

    \emph{Basis.} The paper reports MDS yield and complexity but does not
    close the loop to survival (LQ $\alpha/\beta$ or RBE). Repair models
    (MEDRAS, LEM, MK, PIDE-derived) map DSBs (or MDS) to survival via
    misrepair kinetics. Without that mapping, the paper's central
    quantities cannot be used to design or verify a clinical RBE model.

    \emph{Next steps.} Feed the P2 (LET, MDS, mean $C$, complex-fraction)
    sweep into MEDRAS / LEM-IV / MK repair kinetics and predict
    $\alpha/\beta$(LET) for p and He; compare to the PIDE experimental
    database (Friedrich \emph{et al.}) and to the sibling BNCT-slot
    MedRAS\_analytic outputs. Identify whether MDS complexity beats
    DSB-only count as a survival predictor at matched dose.

  \item \textbf{How sensitive are the MGM per-Gy MDS yields to chromatin
    geometry (heterochromatin fraction, nucleus shape, sub-nuclear domain
    organisation), which the current condensed-history MGM ignores by
    using a spherical uniform 9.65 $\mu$m nucleus?}

    \emph{Basis.} The MGM paper uses a spherical, uniform-DNA-density
    9.65 $\mu$m nucleus; real cells have prolate/oblate shapes, condensed
    heterochromatin fractions of 10--40\%, and non-uniform DNA density.
    Track-structure-aware DSB models (e.g.\ GD or nano-dosimetric) show
    10--30\% sensitivity to these geometric factors. Whether MGM's MDS
    yield is similarly sensitive determines whether cell-line-specific
    corrections are needed for clinical use.

    \emph{Next steps.} (1) Run TOPAS-MGM (once released) with 3 nucleus
    geometries (sphere vs prolate $6 \times 12\,\mu$m vs sphere with 30\%
    dense subvolume) at 3 LETs (10, 30, 60 keV/$\mu$m) for p and He;
    report per-Gy MDS spread. (2) If spread $> 10\%$, publish
    geometry-corrected MDS tables. (3) Cross-check against
    Friedland/Kundrat 2020 DNA-target geometry sweeps.

  \item \textbf{How does MGM compare against other track-structure
    DNA-damage codes (PARTRAC, KURBUC, Geant4-DNA-option-2/4 direct,
    RITRACKS) at matched conditions, and is any code closer to a
    hypothetical FLASH / ultra-high-dose-rate regime where transient
    chemistry is truncated?}

    \emph{Basis.} The paper compares TOPAS-MGM only against TOPAS-nBio
    Geant4-DNA option-2. Cross-code comparison against PARTRAC
    (Friedland), KURBUC (Nikjoo), RITRACKS (Plante) is a standard
    benchmark that the paper skips. Also: FLASH ($>40$ Gy/s) modifies
    transient chemistry (reduced OH-mediated damage); MGM does not track
    dose-rate; whether MGM breaks or holds at FLASH regime is unknown.

    \emph{Next steps.} (1) Assemble a cross-code MDS benchmark at 5 LETs
    (2, 10, 30, 60, 100 keV/$\mu$m) for proton and helium: MGM
    (analytical) + TOPAS-MGM (once released) + PARTRAC + KURBUC +
    Geant4-DNA option-4 direct. (2) Add a dose-rate axis at 3 rates
    (0.1, 1, 100 Gy/s) using Geant4-DNA-Chem stage; test whether MGM's
    LET-only kernel needs a dose-rate correction. (3) Correlate the
    (code, LET, dose-rate) MDS spread against clinical FLASH-sparing
    measurements (Vozenin, Favaudon).
\end{enumerate}
